% poster.tex — A0 portrait poster (tikzposter)
%
% Backup poster for the NeurIPS 2026 Evaluations & Datasets submission.
% Compile with: xelatex poster.tex   (or pdflatex)
%
% Re-formattable to A0 landscape by switching \documentclass orientation.
% Re-formattable to 36" x 48" landscape by setting [size=custom,...].

\documentclass[a0paper,portrait,fontscale=0.32]{tikzposter}

\usepackage[utf8]{inputenc}
\usepackage{lmodern}
\usepackage{tikz}
\usepackage{booktabs}
\usepackage{graphicx}
\usepackage{xcolor}
\usepackage{hyperref}

\usetheme{Simple}
\usecolorstyle{Default}

\title{When Workflow Conditioning Helps Remaining Surgery Duration: MultiBypass140 vs Cholec80, with a Causal-at-Inference Variant}
\author{Bill Chen \and [Co-authors TBD]}
\institute{Xi Lab, [Affiliation]}
\titlegraphic{}

\begin{document}

\maketitle

% ---------- ROW 1: motivation, datasets, method, headline result ---------
\begin{columns}

  \column{0.25}
  \block{1.\ The problem}{
    \textbf{Remaining surgery duration} (RSD) prediction matters
    operationally: OR time costs \$40--100/min.

    \textbf{Workflow heterogeneity} is a hidden driver of error: two
    surgeons performing the same procedure can have different phase
    orderings, pacing, and center habits. A single non-conditioned
    predictor blends regimes and is wrong on non-canonical workflows.

    \textbf{Question.} \emph{When does giving the model an explicit
    workflow signal help, and can we extract it causally at inference?}
  }

  \column{0.25}
  \block{2.\ Two contrasting datasets}{
    \begin{tabular}{lll}
      \toprule
      & \textbf{MB140} & \textbf{Cholec80} \\
      \midrule
      Procedure & RYGB & Cholecystectomy \\
      Centers & 2 & 1 \\
      Videos & 140 & 72 (public) \\
      Median & 83 min & 35 min \\
      Clusters & 6 balanced & 4, one $=$71\% \\
      \bottomrule
    \end{tabular}
    \vspace{0.5em}

    \textbf{MB140 has high workflow variability; Cholec80 has low.}
    That contrast is the test bed.
  }

  \column{0.25}
  \block{3.\ Method}{
    Frozen-lower ViT-B/16 + Surgformer HTA + 3 heads (RSD, phase,
    deviation). Workflow-cluster embedding prepended to frame sequence.

    \textbf{Two ways to compute the cluster:}
    \begin{itemize}
      \item \textbf{Oracle (retrospective)}: from full-video phase
            sequence. Privileged information; not deployable; an upper
            bound.
      \item \textbf{Causal-at-inference}: model's own phase head over
            the prefix $\rightarrow$ TF-IDF + PCA + k-means $\rightarrow$
            soft cluster posterior $\rightarrow$ mixture embedding.
            Only pixels at test time.
    \end{itemize}
  }

  \column{0.25}
  \block{4.\ Headline result (MB140 fold 0)}{
    \begin{tabular}{lr}
      \toprule
      Configuration & Val MAE \\
      \midrule
      No-token & 13.55 \\
      Oracle & 12.59 $\pm$ 0.33 \\
      \textbf{Causal} & \textbf{12.56 $\pm$ 0.04} \\
      \bottomrule
    \end{tabular}
    \vspace{0.5em}

    \textbf{Causal matches oracle, with 8$\times$ tighter cross-seed
    variance.} Paired Wilcoxon: oracle and causal each beat no-token at
    $p=0.026$; causal vs.\ oracle indistinguishable ($p=0.85$).
  }

\end{columns}

% ---------- ROW 2: variability scaling + cross-center -------------------
\begin{columns}

  \column{0.5}
  \block{5.\ The variability-scaling claim}{
    \begin{tabular}{lll}
      \toprule
      Dataset & Workflow diversity & Token effect \\
      \midrule
      MB140 & High (largest cluster 27\%) & $-$0.96 min \\
      Cholec80 & Low (largest cluster 71\%) & $-$0.12 min (null) \\
      \bottomrule
    \end{tabular}
    \vspace{0.5em}

    \textbf{The benefit of workflow conditioning scales with the
    workflow variability of the dataset.} The Cholec80 null is
    \emph{predicted} by the cluster-size histogram alone --- it
    confirms the scaling hypothesis.
  }

  \column{0.5}
  \block{6.\ Cross-center is still hard}{
    \begin{tabular}{lr}
      \toprule
      Bern $\rightarrow$ Strasbourg & Val MAE \\
      \midrule
      No-token & 18.26 \\
      Oracle & 18.05 \\
      Causal & TBD (in progress) \\
      \bottomrule
    \end{tabular}
    \vspace{0.5em}

    Workflow conditioning closes $\sim$0.2 of a 6-minute gap.
    \textbf{Multi-center distributional shift, not workflow alone, is
    the dominant unsolved problem.}
  }

\end{columns}

% ---------- ROW 3: causal pipeline figure -------------------------------
\block{7.\ Causal-at-inference pipeline}{
  \centering
  \begin{tikzpicture}[node distance=1.5cm, every node/.style={align=center}]
    \node (frames)  [draw, rounded corners] {Prefix frames\\$x_{1{:}t}$};
    \node (vit)     [draw, rounded corners, right=of frames]   {ViT-B + HTA};
    \node (phase)   [draw, rounded corners, right=of vit]      {Phase head\\($\hat\phi_{1{:}t}$)};
    \node (bigram)  [draw, rounded corners, right=of phase]    {TF-IDF\\over phase bigrams};
    \node (pca)     [draw, rounded corners, right=of bigram]   {PCA(16) +\\k-means centroids};
    \node (post)    [draw, rounded corners, right=of pca]      {Soft posterior\\$q(z|x_{1{:}t})$};
    \node (mix)     [draw, rounded corners, right=of post]     {$\sum_k q_k \mathbf{E}[k]$\\workflow token};

    \draw[->] (frames) -- (vit);
    \draw[->] (vit)    -- (phase);
    \draw[->] (phase)  -- (bigram);
    \draw[->] (bigram) -- (pca);
    \draw[->] (pca)    -- (post);
    \draw[->] (post)   -- (mix);
  \end{tikzpicture}

  \vspace{0.5em}
  No future frames or ground-truth phase labels at inference. Phase-head
  accuracy 0.78 $\rightarrow$ cluster agreement with oracle 91\%
  $\rightarrow$ RSD MAE matches oracle.
}

% ---------- ROW 4: take-aways, limitations, references ------------------
\begin{columns}

  \column{0.34}
  \block{8.\ Take-aways}{
    \begin{enumerate}
      \item \textbf{Benchmark choice matters}: Cholec80 saturates,
            MB140 exposes the workflow signal.
      \item \textbf{Causal-at-inference is feasible}: the model's own
            phase head recovers the oracle-quality signal.
      \item \textbf{Cross-center is the next problem}: workflow
            conditioning helps only a little; modeling visual-domain
            and operative-style shift is required.
    \end{enumerate}
  }

  \column{0.33}
  \block{9.\ Limitations}{
    \begin{itemize}
      \item Causal at inference, not at training (teacher forcing).
      \item 5-fold CV partially complete; final manuscript reports
            cross-fold mean $\pm$ std.
      \item Cholec80 on 72 of 80 videos (8 lack public phase labels);
            comparisons to TransLocal / Bose are directional, not
            benchmark-clean.
      \item Deviation head $\sim$0.40 F1 (auxiliary, not competitive).
    \end{itemize}
  }

  \column{0.33}
  \block{10.\ Code \& citations}{
    \textbf{Code:} \texttt{brsd\_lib} clean-room library
    (\texttt{causal\_cluster}, \texttt{smoothing}, \texttt{stats},
    \dots).

    \textbf{Anchor citations:}
    Twinanda~et~al.~(2018) RSDNet;
    Loukas~et~al.~(2024) TransLocal;
    Yang~et~al.~(2024) Surgformer;
    Lavanchy~et~al.~(2024) MultiBypass140;
    Vapnik~\&~Vashist~(2009) LUPI.

    \vspace{0.5em}
    \textbf{Full paper}: \texttt{paper/review\_manuscript.docx}
    (anonymized version provided to OpenReview).

    [QR code: paper PDF + code repo]
  }

\end{columns}

\end{document}
